\section*{Guide to the Supplemental Material}

This supplement collects the proofs, estimation details, experimental protocols,
and additional results supporting \emph{Fisher-Glass Learning: Classification
Bounds and Subpopulation Audits}. Appendix~\ref{app:core-proofs} gives the
theoretical arguments, including the worked classification family.
Appendices~\ref{app:estimation-certification} and~\ref{app:hidden-slice-discovery}
define the operational matrices, uncertainty calculations, and frozen-gate
protocol. Appendices~\ref{app:experimental-details}
and~\ref{app:information-interventions} report the empirical checks and
information-changing interventions. The map below locates the completed
comparisons and their principal interpretation limits. Appendix letters and result numbering A--E match the accompanying paper. References to main-text sections and results link to \path{main.pdf}; appendix references are internal. Appendix~\ref{app:supplement-reproducibility} gives the evidence and replay index.

\begingroup
\small
\setlength{\tabcolsep}{4pt}
\renewcommand{\arraystretch}{1.12}
\begin{longtable}{@{}>{\raggedright\arraybackslash}p{.20\linewidth}
                  >{\raggedright\arraybackslash}p{.52\linewidth}
                  >{\raggedright\arraybackslash}p{.23\linewidth}@{}}
\toprule
Comparison & Protocol and result to read together & Location \\
\midrule
\endfirsthead
\toprule
Comparison & Protocol and result to read together & Location \\
\midrule
\endhead
Frozen background intervention &
Validation selects the waterbird/water group before held-out scoring. Same-type
and opposite-type recompositions share bird foregrounds and labels. The
controlled accuracy difference is 52.6 percentage points across five training
seeds; confident-flip results are reported separately. &
Appendix~\ref{app:prospective-intervention};
Table~\ref{tab:d-intervention-outcomes} \\
\addlinespace
Selection and summary controls &
Accuracy and negative log-likelihood (NLL) selectors agree with the surviving-fraction rule on the same
eligible set. Centering reverses the selection; the specified predictive-Fisher
construction preserves it. The four-summary comparison does not establish
superior swap-damage prediction by inverse-CVaR. &
Tables~\ref{tab:d-selection-comparison}--\ref{tab:d-matrix-per-seed},
\ref{tab:four-way-swap} \\
\addlinespace
Auxiliary failure prediction &
Species-held-out prediction adds background and residual coordinates to source
loss, confidence, and logit predictors. Gains within the selected group are
positive across seeds; pooled average precision decreases when group indicators are
included. This does not isolate a scalar-information or residual-only effect. &
Appendix~\ref{app:prospective-intervention};
Tables~\ref{tab:d-flip-prediction}, \ref{tab:supplement-flip-intervals} \\
\addlinespace
Ridge and exact-information audits &
The 75 default-ridge audits report nuisance conditioning, ridge inflation, exact
plug-in values, and nominal conditional Wald lower values. Fresh training
audits are separate from historical training cohorts; numerical lower/exact
ordering alone does not establish confidence coverage. &
Appendix~\ref{app:implemented-audit};
Tables~\ref{tab:appB-conditioning}, \ref{tab:appB-training} \\
\addlinespace
Hidden-slice discovery &
Task labels and a fixed nuisance direction define the search; group annotations
remain sealed through evaluation. Matched-mass selectors, random-direction and
feature-moment controls, support, stability, and post hoc composition are
reported. CLIP versus random direction remains inconclusive. &
Appendices~\ref{app:hidden-slice-discovery}, \ref{app:discovery-results};
Table~\ref{tab:d-hidden-selectors} \\
\addlinespace
Constrained slice repair &
Loss, Fisher, and combined repairs share accuracy constraints and an unchanged
candidate. Fisher repair selects that unchanged model in all five seeds.
Group-supervised Group DRO improves worst-group accuracy without raising the
rediscovered slice's mean information. &
Appendix~\ref{app:extended-training};
Table~\ref{tab:e-repair-performance} \\
\addlinespace
Group-free coefficient selection &
This separate post hoc audit selects the same coefficients as validation-NLL
minimization. Information increases while worst-group accuracy decreases by
12.6 points on average. The proposed threshold-and-risk-budget policy is
distinct and has not been prospectively evaluated. &
Appendix~\ref{app:extended-training};
Table~\ref{tab:e-groupfree-performance} \\
\bottomrule
\end{longtable}
\endgroup

\paragraph{Cohorts and scope.}
The primary prospective intervention uses seeds 22--26 and the surviving-fraction
rule. The earlier development cohort uses seeds 17--21 and an absolute-information
rule; both use the same benchmark. Matched selectors, matrix constructions, and
the four-summary comparison analyze retained records after the intervention
outcomes were known. Hidden-slice discovery and its constrained repair use the
corrected v2 cohort, distinct from the post hoc group-free coefficient cohort.
No retained experiment tests whether the learned low-information slice is more
sensitive than a confidence-matched population. Intervals across seeds describe
fitting or partition variation on fixed data, not independent datasets.

\paragraph{Provenance and replay.}
The accompanying artifact's \path{START_HERE.md} and
\path{docs/MANUSCRIPT_MAP.md} map these results to retained inputs, hashes, and
replay commands. Numerical replay uses saved records and does not retrain the
image classifiers. Historical synthetic, diagnostic, training, and acquisition
summaries have explicitly documented gaps in their original per-seed outputs;
fresh audits and regenerated figures do not replace those missing records.
Acquisition means therefore carry no reconstructed historical intervals.
The experiment guides distinguish independent replication from exact replay.
